\section{Conclusion}
\label{sec:conclusion}

LLM agents are being connected to real HPC schedulers, filesystems, and scientific workflows now, and HPC sites are already noticing and worrying about it \cite{Ma2025HPCAgent,Rosendo2025ORNL,Mondesire2025PEARC,AaltoSciComp2026}. The security question this raises --- what happens when a correctly authenticated, correctly authorized agent is redirected by adversarial content it encounters during normal operation --- is not answered by either of the literatures adjacent to it. LLM-agent security research has the mechanisms but not the HPC-specific surfaces or consequences; HPC security research has the access-control rigor but not a way to represent task-scoped intent beneath account-level authority. We have defined the hijacked authorized agent threat model to name this gap precisely, taxonomized the HPC-specific attack surfaces that follow from it, shown why existing controls do not close it, and outlined a research agenda --- including an empirical benchmark now in development --- to study and eventually mitigate it. The goal is not to declare HPC agents unsafe; it is to make sure the field can tell the difference between an agent that is safe and one that has simply not yet been tested against its own environment.

\subsection*{AI Disclosure}
This paper was prepared with the assistance of AI-powered academic writing and research tools, including literature identification, citation verification, and drafting support. All content, arguments, and conclusions were directed and reviewed by the author(s), who take full responsibility for the accuracy and integrity of this work.
